\documentclass[12pt,a4paper]{article}
\usepackage[utf8]{inputenc}
\usepackage{amsmath,amssymb,amsthm}
\usepackage{algorithm}
\usepackage{algpseudocode}
\usepackage{geometry}
\usepackage{hyperref}
\usepackage{booktabs}

\geometry{margin=2.5cm}

\newtheorem{theorem}{Theorem}
\newtheorem{lemma}[theorem]{Lemma}
\newtheorem{corollary}[theorem]{Corollary}
\newtheorem{definition}{Definition}

\title{Problem 451: Modular Inverses}
\author{}
\date{}

\begin{document}
\maketitle

\section{Problem Statement}

For a positive integer $n$, define $I(n)$ as the largest positive integer $m < n-1$ such that
$m^2 \equiv 1 \pmod{n}$. If no such $m$ exists, set $I(n) = 1$.

Known values: $I(7) = 1$, $I(15) = 11$, $I(100) = 51$.

Compute
\[
    \sum_{n=3}^{2 \times 10^7} I(n).
\]

\section{Mathematical Foundation}

\begin{definition}
A number $m$ with $1 \le m < n$ and $\gcd(m, n) = 1$ is a \emph{self-inverse} modulo $n$ if
$m^2 \equiv 1 \pmod{n}$.
\end{definition}

\begin{theorem}[CRT factorization of self-inverses]\label{thm:crt}
Let $n = p_1^{a_1} p_2^{a_2} \cdots p_k^{a_k}$. Then
\[
    m^2 \equiv 1 \pmod{n} \iff m^2 \equiv 1 \pmod{p_i^{a_i}} \text{ for all } i = 1, \ldots, k.
\]
\end{theorem}

\begin{proof}
By the Chinese Remainder Theorem, $\mathbb{Z}/n\mathbb{Z} \cong \prod_{i=1}^{k} \mathbb{Z}/p_i^{a_i}\mathbb{Z}$. The equation $m^2 - 1 \equiv 0$ holds in the product ring if and only if it holds in each factor.
\end{proof}

\begin{lemma}[Solutions modulo prime powers]\label{lem:pp}
The solutions of $x^2 \equiv 1 \pmod{p^a}$ are:
\begin{itemize}
    \item For $p$ odd: exactly 2 solutions, $x \equiv \pm 1 \pmod{p^a}$.
    \item For $p = 2$:
    \begin{itemize}
        \item $a = 1$: 1 solution ($x = 1$).
        \item $a = 2$: 2 solutions ($x \in \{1, 3\}$).
        \item $a \ge 3$: 4 solutions ($x \in \{1,\; 2^{a-1}-1,\; 2^{a-1}+1,\; 2^a-1\}$).
    \end{itemize}
\end{itemize}
\end{lemma}

\begin{proof}
For odd $p$, $(\mathbb{Z}/p^a\mathbb{Z})^*$ is cyclic of order $\varphi(p^a) = p^{a-1}(p-1)$. In a cyclic group, the equation $x^2 = e$ has at most 2 solutions, and $\pm 1$ are always solutions, so there are exactly 2.

For $p = 2$ and $a \ge 3$, we have $(\mathbb{Z}/2^a\mathbb{Z})^* \cong \mathbb{Z}/2 \times \mathbb{Z}/2^{a-2}$, so the equation $x^2 = 1$ has $2 \times 2 = 4$ solutions (the identity and the unique element of order 2 in each cyclic factor). These are verified directly:
\begin{align*}
    1^2 &\equiv 1, \\
    (2^{a-1}-1)^2 &= 2^{2a-2} - 2^a + 1 \equiv 1, \\
    (2^{a-1}+1)^2 &= 2^{2a-2} + 2^a + 1 \equiv 1, \\
    (2^a - 1)^2 &= 2^{2a} - 2^{a+1} + 1 \equiv 1 \pmod{2^a}.
\end{align*}
The cases $a = 1$ and $a = 2$ are checked by exhaustion.
\end{proof}

\begin{corollary}
The total number of self-inverses modulo $n = \prod p_i^{a_i}$ is $\prod_{i=1}^{k} s(p_i^{a_i})$, where $s(p^a)$ denotes the count from Lemma~\ref{lem:pp}.
\end{corollary}

\begin{theorem}[Characterization of $I(n)$]\label{thm:In}
Given the complete set of self-inverses modulo $n$ (constructed via CRT from Theorem~\ref{thm:crt} and Lemma~\ref{lem:pp}), $I(n)$ equals the second-largest element of this set. The largest is always $n-1$; if the only solutions are $\{1, n-1\}$, then $I(n) = 1$.
\end{theorem}

\begin{proof}
Since $(n-1)^2 = n^2 - 2n + 1 \equiv 1 \pmod{n}$, the value $m = n-1$ is always a self-inverse. By definition, $I(n)$ is the largest self-inverse strictly less than $n-1$. If no non-trivial self-inverse exists, the only candidate below $n-1$ is $m = 1$, giving $I(n) = 1$.
\end{proof}

\section{Algorithm}

\begin{algorithm}[H]
\caption{Compute $\sum I(n)$ via SPF sieve and CRT}
\begin{algorithmic}[1]
\State Compute smallest prime factor $\text{spf}[i]$ for $2 \le i \le N$ via sieve of Eratosthenes.
\State $\text{total} \gets 0$
\For{$n = 3$ to $N$}
    \State Factor $n$ using $\text{spf}$: obtain prime powers $p_1^{a_1}, \ldots, p_k^{a_k}$.
    \State $\text{sols} \gets \{1\}$, $\text{mod} \gets 1$
    \For{each $(p_i, p_i^{a_i})$}
        \State Compute local solutions $L_i$ of $x^2 \equiv 1 \pmod{p_i^{a_i}}$ via Lemma~\ref{lem:pp}.
        \State $\text{sols} \gets \{\ \text{CRT}(s_1, \text{mod},\ s_2, p_i^{a_i}) : s_1 \in \text{sols},\ s_2 \in L_i\ \}$
        \State $\text{mod} \gets \text{mod} \times p_i^{a_i}$
    \EndFor
    \State $I(n) \gets \max\{s \in \text{sols} : 1 < s < n-1\}$ (or $1$ if no such $s$)
    \State $\text{total} \gets \text{total} + I(n)$
\EndFor
\State \Return total
\end{algorithmic}
\end{algorithm}

\section{Complexity Analysis}

\begin{itemize}
    \item \textbf{Time:} $O(N \log \log N)$ for the SPF sieve. For each $n$, factoring via SPF takes $O(\log n)$. The CRT reconstruction processes $\omega(n)$ prime factors with up to $2^{\omega(n)}$ solutions. Since $\omega(n) \le O(\log n / \log \log n)$ and on average $\omega(n) \approx \log \log n$, the amortized cost per $n$ is small. Overall: $O(N \log \log N)$.
    \item \textbf{Space:} $O(N)$ for the SPF array.
\end{itemize}

\section{Answer}

\[
\boxed{153651073760956}
\]

\end{document}
